\begin{problem}[Gluing along principal affine opens]\label{prob:vi20-p06}
Let $A_f\cong B_g$. Construct the scheme obtained by gluing $\Spec A$ and $\Spec B$ along $D(f)\cong D(g)$.
\end{problem}

\ifdefined\FullProblemDossiers
\paragraph{Why this problem matters.}
This is the algebraic normal form for two-chart gluing.

\paragraph{Useful ideas before starting.}
Use VI/18 to identify the overlap schemes and VI/19 to reverse the ring isomorphism.

\paragraph{Guided hints.}
\begin{enumerate}
\item Use VI/18 to identify the overlap schemes and VI/19 to reverse the ring isomorphism.
\end{enumerate}

% VI20_FULL_SOLUTION_REFINEMENT P06
\begin{solution}
Choose an isomorphism of rings
\[
\theta:B_g\xrightarrow{\sim}A_f.
\]
By the principal-open theorem,
\[
D(f)\subseteq\Spec A\cong\Spec A_f,
\qquad
D(g)\subseteq\Spec B\cong\Spec B_g.
\]
Contravariance of $\Spec$ turns $\theta$ into an isomorphism of affine schemes
\[
\varphi:D(f)\xrightarrow{\sim}D(g).
\]
Use $\varphi$ and $\varphi^{-1}$ as the transition maps in the two-chart gluing theorem.

The resulting space is the quotient of $\Spec A\amalg\Spec B$ in which a prime of $A$ avoiding $f$ is identified with the corresponding prime of $B$ avoiding $g$ under $\Spec(\theta)$.  Both original spectra remain open subschemes.  On an open set meeting both charts, a section is a pair of local sections whose restrictions to the overlap correspond through $\theta$ (and its localizations).  In particular, when the whole overlap is used, the compatibility condition on functions is exactly the equality induced by
\[
B_g\xrightarrow{\theta}A_f.
\]
Thus the localized ring isomorphism supplies the full transition data: points, topology, and regular functions.
\end{solution}


\paragraph{What happened mathematically.}
Localized coordinate rings are transition functions for affine charts.

\paragraph{Extensions.}
Compute global sections as an equalizer.

\paragraph{Provenance.}
\textbf{New canonical synthesis; no numbered legacy Problem is claimed.}
\else
\paragraph{Provenance.} \textbf{New canonical synthesis; no numbered legacy Problem is claimed.}
\fi
